\section{Related work}\label{sec:related}
Prognostic modelling of NSCLC from pretreatment CT has progressed from clinical and radiomic models, through deep networks trained on tumour-centred patches, to foundation models pretrained at scale and adapted to the prognostic task. This section first traces the transition from clinical and radiomic models to deep learning, and the level of discrimination these approaches have reached (Section~\ref{sec:related-repr}). It then reviews CT foundation models and the strategies used to adapt them (Section~\ref{sec:related-fm}), and the evaluation of their prognostic performance, generalization and tumour dependence (Section~\ref{sec:related-eval}). It closes by positioning the present study against the closest prior work (Section~\ref{sec:related-position}).

\subsection{From clinical and radiomic models to deep learning}\label{sec:related-repr}
Stage and tumour size are the established prognostic factors in NSCLC. Stage is assigned with the TNM classification, which combines primary-tumour extent (T), nodal involvement (N) and distant metastasis (M) into stage groups across which survival decreases.\cite{goldstraw2016} In patients treated with radiotherapy, gross tumour volume is also associated with survival.\cite{bradley2002} Other established predictors describe the patient and the treatment rather than the tumour: in NSCLC treated with radiotherapy or chemoradiotherapy, performance status, lung function, the number of positive lymph-node stations on PET, blood biomarkers and the delivered dose contribute to survival models.\cite{dehingoberije2009,dehingoberije2010,oberije2015} Even with these variables, discrimination remains moderate: a model combining several of them, developed at MAASTRO, the institution that provided LUNG1, reached a concordance of 0.62 after correction for optimism and 0.58 to 0.60 in two external cohorts, only modestly above the 0.50 to 0.55 of models based on TNM stage alone.\cite{oberije2015} 

Radiomics extends these measures with handcrafted descriptors of tumour shape, intensity and texture extracted from the segmented tumour. Aerts and colleagues introduced the LUNG1 cohort with a radiomic signature that was prognostic in independent lung and head-and-neck cancer cohorts.\cite{aerts2014} Radiomic features have since remained a strong comparator on LUNG1, and the learned image representations discussed below have not clearly surpassed them.\cite{braghetto2022,neukirch2026} Clinical and radiomic models are therefore the reference against which learned representations need to be assessed.

The first learned representations for CT-based NSCLC prognosis kept the tumour-centred view of radiomics: deep networks have largely operated on patches centred on the primary tumour. Hosny and colleagues trained a three-dimensional convolutional network on radiotherapy and surgical cohorts and tested it on the MAASTRO cohort that forms part of NSCLC-Radiomics.\cite{hosny2018} Pai and colleagues pretrained a lesion encoder contrastively on DeepLesion and evaluated frozen features for two-year overall survival on independent test sets from NSCLC-Radiomics and NSCLC-Radiogenomics.\cite{pai2024} Similar lesion-encoder features have been compared with radiomic and clinical models in a multicentre radiotherapy trial cohort.\cite{upadhaya2024} Both evaluation cohorts in the present study have therefore served as prognostic test sets before. Across these studies, handcrafted features were replaced by learned representations, but the input remained centred on the lesion and survival was mostly reduced to a binary endpoint at a fixed horizon. On LUNG1, these approaches have reached a similar level. A comparison of radiomic features, two-dimensional deep features and convolutional networks for two-year survival found comparable results across pipelines (test AUC 0.63 to 0.67) and concluded that further imaging information is difficult to extract from this cohort.\cite{braghetto2022} Its authors proposed three-dimensional models, contrastive pretraining and analysis of the whole lung, including the tissue around the tumour, as directions for further work.

\subsection{Foundation models for CT-based prognosis}\label{sec:related-fm}
Recently released general-purpose encoders for three-dimensional CT follow these directions: they encode the whole volume and are pretrained on large collections, several of them with contrastive objectives. They differ in their pretraining data and objective. CT-CLIP aligns chest CT volumes with radiology reports from the CT-RATE dataset.\cite{hamamci2026ctclip} COLIPRI extends language--image pretraining with report generation and vision-only objectives, using CT-RATE together with NLST.\cite{wald2025colipri} TANGERINE is a volumetric masked autoencoder pretrained on approximately 98,000 chest CT volumes from lung-screening and other public and private collections.\cite{mcconnell2026tangerine} Rad-JEPA-3D uses a joint-embedding predictive objective on CT volumes from M3D-Cap.\cite{trinh2026radjepa} Beyond these CT-specific encoders, foundation models trained on natural images, such as DINOv3, provide a general visual representation that can also be transferred to CT.\cite{simeoni2025dinov3} Because such encoders operate on two-dimensional images, they are applied to a CT volume slice by slice, and the slice representations are then aggregated; on LUNG1, aggregating them through soft attention outperformed benchmark three-dimensional networks for two-year overall survival.\cite{paolo2026}

Whether volumetric or slice-based, these foundation models are pretrained without survival labels and must therefore be adapted to the prognostic task, which can be done in several ways. The encoder can be kept frozen and a supervised head fitted on its features; low-rank matrices can be added to the attention layers and trained while the original weights stay fixed (LoRA);\cite{hu2021} or part or all of the encoder can be fine-tuned. For prognosis, the available evidence favours the frozen encoder. Pai and colleagues found that frozen lesion-encoder features outperformed fine-tuned variants on the prognostic task.\cite{pai2024} Allgöwer and colleagues used whole-volume CT-CLIP representations with a survival head in a single-centre lung cancer cohort; they compared the frozen encoder, LoRA and full fine-tuning and reported no advantage of adaptation.\cite{allgower2026}

This evidence remains limited. Each study examined a single encoder, and neither examined how the comparison depends on the number of training labels, the optimization schedule or checkpoint selection, or how it transfers to an external cohort. The original evaluations of COLIPRI, TANGERINE and Rad-JEPA-3D concern detection, classification, segmentation, retrieval, report generation and visual question answering; we did not identify a survival evaluation of these encoders. The availability of pretrained encoders therefore does not settle which adaptation strategy works best for survival prediction when labelled data are limited.

\subsection{Prognostic evaluation, generalization and tumour dependence}\label{sec:related-eval}
Two recent preprints benchmark foundation-model features on the cohorts used here. TumorImagingBench compared frozen, lesion-centred embeddings from ten CT foundation models, including CT-FM,\cite{pai2025ctfm} CT-CLIP and the Pai lesion encoder. Its tasks included two-year survival on NSCLC-Radiomics and NSCLC-Radiogenomics with a nearest-neighbour head, and it reported only modest discrimination across models.\cite{aerts2025tib} Neukirch and colleagues developed models on LUNG1 and tested them on an external cohort. They compared Curia, DINOv3 and radiomic features with several classification heads, and two-year survival was one of their tasks.\cite{neukirch2026} Outside the lung, a similar frozen-feature design has been applied to head-and-neck cancer recurrence with external testing.\cite{guetarni2026} These benchmarks use binary endpoints, which discard the follow-up of patients censored before the horizon, rather than survival models that handle censoring explicitly, and evaluate only heads fitted on frozen encoders.

Some of these benchmarks also test their models on an external cohort, which is needed to judge whether a prognostic model transfers beyond its development cohort; differences in case mix between cohorts then determine how its performance should be interpreted.\cite{debray2015,collins2024} Transfer has two components. Discrimination measures whether the model ranks patients correctly by risk; calibration measures whether its predicted absolute risks agree with observed outcomes. A model can retain its ranking in a new cohort while systematically over- or underestimating risk.\cite{vancalster2019} For survival models, calibration is assessed at a fixed time horizon against observed event rates estimated under censoring.\cite{austin2020} A further question is whether an image representation adds prognostic value to the clinical model, which is what determines its potential clinical use. The benchmarks above report discrimination only; we did not identify one that reports calibration or incremental value over clinical models on an external cohort.

Beyond whether predictions transfer, it matters what they are based on: encoders that see the whole CT may draw on the tumour or on the surrounding anatomy, and whether CT-based survival predictions rely on the tumour has been examined with ablation and saliency methods. Hosny and colleagues retrained their network after masking the region outside the tumour, which reduced discrimination, and visualized activation maps concentrated at the tumour boundary.\cite{hosny2018} TumorImagingBench reported occlusion-based saliency that localized to the lesion for some encoders but not others.\cite{aerts2025tib} Neukirch and colleagues compared tumour-mask, lung-mask and unsegmented inputs.\cite{neukirch2026} A frozen-feature benchmark of volumetric CT encoders found that small, low-contrast focal findings are poorly represented after global pooling.\cite{chevli2026} These analyses either retrain the model on modified inputs or perturb the image. None intervenes on the internal representation of a frozen encoder, which would separate the contribution of the tumour from that of the surrounding anatomy without changing the model. In summary, the existing benchmarks do not combine censoring-aware survival analysis, external evaluation with calibration, and an analysis of tumour dependence by token masking inside frozen encoders under matched spatial controls.

\begin{table}[!htbp]
\centering\footnotesize
\caption{Comparison with the closest prior work on foundation-model features for CT-based survival prediction in lung cancer. Entries summarize the published descriptions of each study.}
\label{tab:positioning}
\setlength{\tabcolsep}{2.5pt}
\newcolumntype{L}[1]{>{\raggedright\arraybackslash}p{#1}}
\begin{tabular}{@{}L{1.8cm}L{1.5cm}L{1.85cm}L{1.45cm}L{1.8cm}L{1.4cm}L{1.65cm}L{1.3cm}L{1.6cm}@{}}
\toprule
Study & Input & Models & Endpoint & Adaptation & Label budget & Tumour dependence & External cohort & Calibration \\
\midrule
Pai et al.\cite{pai2024} & Lesion patch & One lesion FM & Binary, 2 years & Frozen; fine-tuning & Not reported$^\dagger$ & No & Yes & Not reported$^\dagger$ \\
Tumor-\newline Imaging-\newline Bench\cite{aerts2025tib} & Lesion patch & Ten CT FMs & Binary, 2 years & Frozen (kNN) & Not reported$^\dagger$ & Occlusion saliency & Not reported$^\dagger$ & No \\
Neukirch et al.\cite{neukirch2026} & Tumour, lung or whole CT & Curia, DINOv3, radiomics & Binary, 2 years & Frozen (7 heads) & No & Input masking & Yes & Not reported$^\dagger$ \\
Allgöwer et al.\cite{allgower2026} & Whole CT & CT-CLIP & Survival$^\dagger$ & Frozen; LoRA; full fine-tuning & No & No & No & Not reported$^\dagger$ \\
Present study & Whole CT; tumour patch or region & Five FMs, two tumour-centred models, radiomics, clinical & Cox, Uno C at 36 months & Frozen and trainable heads; LoRA; partial fine-tuning & 10--100\% & Token masking with matched controls & Yes & Yes \\
\bottomrule
\end{tabular}

\smallskip
\raggedright FM, foundation model; kNN, $k$-nearest neighbours. $^\dagger$To be confirmed against the full text.
\end{table}

\gap{Confirm the table entries marked $\dagger$ against the full texts. Before submission, repeat the literature search, confirm full-text details of the preprints (including the identity of the external cohort in Neukirch et al.), and replace preprint citations with peer-reviewed versions where available. Verify the author lists of oberije2015, dehingoberije2009, dehingoberije2010 and braghetto2022 against the publishers' records; dehingoberije2009 and dehingoberije2010 were checked against the abstracts only.}

\subsection{Positioning of the present study}\label{sec:related-position}
Prior work on CT-based survival prediction in NSCLC has thus evaluated each family of models in its own setting, with different cohorts, inputs, endpoints and baselines, so the families cannot be compared directly. Table~\ref{tab:positioning} compares the present study with the closest prior work. Each previous study addresses some of the components considered here, but none combines them. The contribution of the present study lies in their joint evaluation within one protocol: several whole-volume and general-purpose foundation models alongside tumour-centred, radiomic and clinical comparators, frozen and adapted models across training-label budgets, censoring-aware survival analysis, tumour-token masking inside the encoders, and external discrimination, incremental clinical value and calibration. Because several of the questions concern whether adaptation or the tumour tokens bring any benefit, equivalence is tested against a prespecified margin, so that the absence of a benefit can be established directly rather than inferred from a nonsignificant difference; none of the studies above reported such tests. Finally, the joint evaluation tests whether the plateau observed on LUNG1 for clinical, radiomic and two-dimensional deep models persists when the directions proposed to overcome it, three-dimensional encoders pretrained at scale on the whole CT, are applied under the same protocol. If it does, the limit to prediction would lie in the information available in the input rather than in the model.
